\subsection{Scholarly and Researcher Recommendation}
Scholarly recommendation has covered papers, experts, collaborators, reviewers, venues, and research groups, combining researcher profiles with multilevel information mining~\cite{yang2014multilevel,sun2015raf} and grouping methods into content-based, collaborative, graph-based, and hybrid approaches~\cite{zhang2023scholarly,bai2019scientific,beel2016research,marcuzzo2022recommendation}. Interaction-based methods represent collaboration patterns but are less useful when historical links are unavailable; publication-based profiling avoids this, as in the Semantic Academic Profiler~\cite{viegas2022semantic}. These studies support publication-based representation, but almost none directly compare several topic models for ranking research groups under one pipeline -- a gap the comparative literature reviewed next has not filled, having instead concentrated on domains outside scholarly recommendation.

\subsection{Topic Modeling for Scholarly Text}
Topic modeling provides one family of content representations for scholarly text: LDA remains a common probabilistic model for latent themes~\cite{jelodar2019lda,blei2003lda}, NMF represents documents through non-negative matrix decomposition~\cite{lee2000nmf}, and BERTopic combines contextual embeddings beyond word co-occurrence~\cite{vaswani2017attention} with UMAP, HDBSCAN, and class-based TF-IDF~\cite{grootendorst2022bertopic}. Domain-adapted embeddings such as SciBERT~\cite{beltagy2019scibert} and the citation-informed SPECTER embeddings used here~\cite{cohan2020specter} better capture scientific-publication semantics than general-domain embeddings, motivating BERTopic's SPECTER-based encoder.

Comparative studies have evaluated LDA, NMF, BERTopic, and related models on short, non-scholarly text such as Twitter posts, online communities, and news headlines, using topic diversity, coherence, or classification accuracy~\cite{egger2022topic,kaur2024moving,babalola2024comprehensive,laureate2023systematic}. Closer to the present domain, Rumadi et al.~\cite{rumadi2025coherence} compared topic coherence across LDA, NMF, BERTopic, and Top2Vec on scientific literature, extending this literature to scholarly text; even this closest prior work, however, evaluates only intrinsic coherence and not a downstream ranking task. Together, this literature shows that model behavior depends on the corpus and evaluation objective without establishing which representation is most useful once publication vectors are aggregated into researcher and research-group profiles. Content-based recommender systems have long compared mean-pooled item or user vectors by cosine similarity to rank candidates~\cite{lops2011content} -- the same mechanism adopted here -- but it has not previously been evaluated with topic vectors from probabilistic, matrix-factorization, and contextual models side by side.

\subsection{Topic Coherence and Downstream Evaluation}
Topic coherence has often been used to select and compare topic models; measures such as $C_v$ estimate whether leading topic words occur in a semantically consistent context~\cite{roder2015exploring,mimno2011optimizing}, though strong intrinsic quality does not necessarily predict downstream performance~\cite{rudiger2022topic}. Recommendation evaluation instead measures whether relevant items appear near the top of a ranked list, using Precision, Recall, F1-score, NDCG, and MRR~\cite{valcarce2020assessing}. These have largely been studied as separate concerns; Section~\ref{sec:results} revisits the separation by reporting both $C_v$ coherence and 13 recommendation metrics for the same three models.
